\chapter{The Non-Aligned: How the Rest of the World Chooses}
\chaptermark{The Non-Aligned}
\label{ch:nonaligned}

Every chapter so far has described a contest between two principals. This one argues that the outcome will be settled by everybody else. The industrial history of Part~I is unambiguous on the point: solar was decided in German feed-in tariffs and then in desert tenders; batteries in the procurement rules of a dozen national auto markets; electric vehicles in Southeast Asia, Latin America, and the European price war. In each case the producer with the best cost curve won because \emph{third markets bought it}, and the incumbent's home market was too small to matter by the time it noticed. AI will be decided the same way, and this chapter therefore treats the non-aligned not as spectators but as the electorate.

The framing correction matters for the book's own recommendations. The strategic question for a country that is neither the United States nor China is not ``which stack do we join?'' It is: \emph{can we evaluate, fork, certify, and operate components from both, while retaining sovereign scientific and industrial capability?} That is a capability question, not an allegiance question, and it has a concrete answer for each actor.

\section{The decision matrix}

Any serious national choice runs through eight dimensions. Stating them explicitly does two things: it exposes the trade-offs that bloc-based framing hides, and it makes different countries' choices comparable.

\begin{table}[htbp]
  \centering
  \small
  \caption{The eight dimensions of a national stack decision. The right-hand column names the chapter of this book that supplies the evidence.}
  \label{tab:decision}
  \begin{tabular}{p{28mm}p{86mm}p{24mm}}
    \toprule
    Dimension & Question & Evidence \\
    \midrule
    Capability & Does the stack satisfy our highest-value workloads---not the benchmark, the workload? & Ch.~\ref{ch:model}, \ref{ch:prize} \\
    Cost & What is the full cost per \emph{successful task}, including operation, integration, and retries? & \S\ref{sec:cost-solved} \\
    Sovereignty & Can weights, runtime, data, and hardware be operated independently of any foreign decision? & Ch.~\ref{ch:compute}, \ref{ch:sizing} \\
    Trust & Can provenance, integrity, behaviour, and liability be certified---by us or by someone we trust? & Ch.~\ref{ch:trust}, \ref{ch:order66} \\
    Interoperability & Does it integrate with our existing cloud, scientific, and enterprise systems? & Ch.~\ref{ch:compute} \\
    Export control & Is continued access exposed to unilateral restriction by a foreign government? & Ch.~\ref{ch:semi}, \ref{ch:governance} \\
    Standards & Does adoption create long-term dependence on one ecosystem's interfaces? & Ch.~\ref{ch:governance} \\
    Industrial spillover & Does adoption build domestic capability, or merely import a service? & Ch.~\ref{ch:theory} \\
    \bottomrule
  \end{tabular}
\end{table}

Two dimensions deserve emphasis because they are routinely omitted from national AI strategies. \emph{Export-control exposure} is asymmetric across the stacks and zero for neither, and Section~\ref{sec:two-commons} prices it row by row. \emph{Industrial spillover} is the dimension that separates a strategy from a purchase: importing tokens builds nothing, while operating models, running probes, and contributing to standards builds people, and people are the only input that compounds.

\section{Two commons, not one}
\label{sec:two-commons}

Everything above was written for a world in which the menu had two items with different shapes: a Chinese commons, free to download and cheap to serve but unaudited and jurisdictionally awkward, and a Western product, expensive and metered but certified and integrated. Chapter~\ref{ch:mirror} ends that world. There are now two commons, and the second one arrives with its training data attached and a chip vendor's balance sheet behind it. NVIDIA's Nemotron family ships open weights together with the training datasets and the reinforcement-learning libraries, under a permissive licence, through the same hubs that carry the Chinese releases; AMD's Instella line is fully open in the strict sense---weights, data and training code; the accompanying open pretraining collection exceeds ten trillion tokens~\cite{nemotron,amd-open} [P]. For a ministry choosing a stack next quarter this is a different menu, and the eight dimensions of Table~\ref{tab:decision} have to be re-run against it rather than extended by one column.

Run them, and the two options resolve into a trade rather than a preference. \emph{Nemotron with open data and CUDA underneath} scores unusually well on trust, for a reason that has nothing to do with benchmarks: a model released with its corpus and its recipes is auditable and, in principle, re-derivable, which a weights-only release is not. A buyer who can retrain is a buyer who can inspect, and this is the first time that dimension has moved toward a \emph{Western} artifact rather than away from one. It scores badly on export-control exposure and on industrial spillover, and the two failures are one failure: the model is free because the complement is not, and the complement is the most tightly controlled product in world trade. The buyer acquires a permanent accelerator dependency priced by a single vendor at a seventy-five per cent gross margin, whose availability to that buyer is set in Washington and revised without notice. \emph{GLM, Qwen, Kimi and their peers with open weights on a capacity-first hardware path} score better on cost, on operational sovereignty and on export-control exposure---a weight file already downloaded cannot be revoked by anyone, and the hardware underneath it is being deliberately designed for parts that are not embargoed (Chapter~\ref{ch:compute}). They score worse on data provenance, since the weights ship without the corpus; worse on trust, for the reasons Chapters~\ref{ch:trust} and~\ref{ch:order66} document; and worse on the question a domestic regulator asks first, which concerns a jurisdiction rather than a model. Neither column dominates, which is better news for the buyer than either column winning would be.

But the second commons has a property the first does not, and a procurement officer who misses it will misprice the option. \emph{It has a single principal contributor with a direct commercial interest in the complement.} Openness supplied by the vendor of the complement is openness with a direction: it is free at the model layer precisely because it is not free at the silicon layer, and the sponsor has said so on the record, repeatedly and without embarrassment---nearly all open models run on its hardware, and its position in open models is good because its ecosystem is everywhere~\cite{nvda-open-models} [P]. A buyer taking Nemotron is therefore not taking a neutral good. It is accepting a subsidy whose purpose is to make a different purchase inevitable, which is the mechanism of Section~\ref{sec:io} operating exactly as the theory predicts and not a scandal. What follows is practical rather than moral: the durability of this openness is a function of the sponsor's strategy rather than of the licence, and the place to monitor it is the sponsor's earnings call. Two things cut the other way and belong in the same breath. A permissive licence on an artifact already released cannot be withdrawn, so the risk is never the loss of what exists but the discontinuation of what comes next---a risk a buyer neutralizes cheaply by mirroring what it depends on. And the same risk applies to the Chinese commons, whose sponsors operate under an industrial policy that could restrict release---though its supply is at least plural, three of the four significant open-weight releases in the last week of August being Chinese and from different firms~\cite{hunyuan-hy4} [V]. A commons with many sponsors is more robust than a commons with one, whichever flag is on it.

The reported acquisition of Hugging Face is the open item, and it belongs in this chapter as a governance fact rather than a commercial one. \emph{The Information} reported on 26~August 2026 that NVIDIA was closing on a purchase of the hub for approximately \$12.9B; at \datafreeze\ neither company had confirmed it, and other reporting held that no agreement had been signed and that the transaction could still fail~\cite{nvda-hf} [A]. This book does not assert that it happened, and the guidance below is written to be correct either way. If it confirms, the distribution point through which very nearly all open weights reach very nearly all users would belong to the vendor of the complement---and what would change is not the commercial terms, which would stay free, but the default discovery surface: what appears on the front page, which artifact is easy to find, which model card is complete, whether a given set of weights remains hosted at all, decided by a party that earns on one of the answers. For a ministry that is a dependency of the same class as a single-vendor cloud region, and it is mitigated the same way and just as cheaply: hold your own copies of the weights, the datasets and the evaluation harnesses you depend on, in a repository you operate, whether or not the report ever confirms.

Which yields an instruction the decision matrix alone does not. A procurement questionnaire for an open model asks about the licence, which is the least informative thing about it. It should also ask about the \emph{provenance of the openness}: who funded the training run, and at which layer do they earn; is the training data released, so that the artifact can be re-derived rather than merely fine-tuned; what happens to the \emph{next} release if the sponsor's strategy changes, and what is the fallback if it stops; is the distribution channel independent of both the sponsor and the vendor of the complement, and do we hold a mirror that survives its loss; and is the model's evaluation performed by anybody who does not earn on the answer. The last is the question no vendor in either commons can answer, and it is why the certification vacancy this chapter keeps returning to has become more valuable rather than less: two competing commons increase the supply of free artifacts and do nothing whatever to increase the supply of independent assurance about them.

\begin{table}[htbp]
  \centering\footnotesize
  \caption{The two commons, as a buyer must price them at \datafreeze. Neither column dominates; the rows are where the difference lives.}
  \label{tab:two-commons}
  \begin{tabular}{p{28mm}p{50mm}p{50mm}}
    \toprule
    Question & Vendor-sponsored Western open stack & Chinese open stack \\
    \midrule
    What ships & Weights, training data, recipes; permissive licence & Weights; licence varies; corpus withheld \\
    \addlinespace
    Who funds it, and where they earn & One accelerator vendor and a rival; on the silicon underneath & Several labs and sponsors; on hardware volume, cloud and standards, later \\
    \addlinespace
    Hardware implied & CUDA and the controlled complement & Capacity-first, deliberately unembargoed \\
    \addlinespace
    Export-control exposure & High---on the complement, not the model & Low on the model; the file cannot be recalled \\
    \addlinespace
    Continuity risk & Sponsor's strategy; one balance sheet & Industrial policy; but plural suppliers \\
    \addlinespace
    Distribution & Hubs; ownership an open question~\cite{nvda-hf} & Same hubs, plus domestic mirrors \\
    \addlinespace
    Independent assurance & None supplied & None supplied \\
    \bottomrule
  \end{tabular}
\end{table}

\section{The export offer, as it actually exists}
\label{sec:export-offer}

Everything to this point has described the non-aligned buyer choosing between artifacts. The story told about that buyer in 2026---in Washington, in Shanghai, and in most of the commentary between them---is that the choice is no longer between artifacts but between \emph{stacks}: a financed, packaged, sovereign-branded Chinese offer of chips, datacentres, models and training, against an American one, with the telecom campaign of Chapter~\ref{ch:telecom} as the template Beijing is supposed to be replaying. This section tests that story against what had actually been signed at \datafreeze, and it does so in the order the record demands: the institutional wrapper first, then the live cases, then the two offers side by side, and then the three points at which the telecom template does and does not transfer. The finding is stated at the end and it cuts in both directions.

\subsection{The wrapper}

The institutional apparatus is real and it is three documents deep. The \emph{Global AI Governance Action Plan} was issued at WAIC on 26~July 2025; the \emph{AI+ International Cooperation Initiative} followed on 23~September 2025; and on 16~July 2026, in Shanghai, with Wang Yi signing for China and the UN Secretary-General in attendance, twenty-nine states signed the founding agreement of the World Artificial Intelligence Cooperation Organization---an intergovernmental body headquartered in Shanghai with a Council of members and a Secretariat, confined by its charter to the civilian domain, funded by member contributions scaled to Universal Postal Union classes, entering into force on three ratifications, and leavable on a year's notice~\cite{waico} [P]. The signatories are reported to include Russia, Brazil, South Africa, Indonesia, Kazakhstan, Pakistan, Malaysia, Serbia, Kenya and most of ASEAN; no official roster was retrievable, and the absences are as informative as the roster: India, the Gulf states other than Oman, Turkey and Egypt did not sign~\cite{waico} [A]. What none of the three texts contains is the thing the sovereign-pact story requires. There is no compute commitment, no financing clause, no datacentre programme and no open-weights undertaking in the Action Plan, in the Initiative, or in the founding agreement. Xi Jinping's WAIC keynote of 17~July pledged 5{,}000 training places for developing-country personnel over five years, regional ``application cooperation centres,'' and a meteorological early-warning system for thirty countries---and nothing else~\cite{waico} [P]. An academic mapping of fifteen comparable international bodies scores WAICO lowest of all of them on formalization and describes it, at founding, as ``a name, a host city, and a statement of purpose''~\cite{waico} [A]. That is not a criticism this book needs to sharpen. The Universal Postal Union dues scale is the tell: it is the funding formula of an organization designed to exist rather than to spend.

\subsection{The three live cases, graded}

Three cases were live at \datafreeze, and they are graded separately because they are not the same kind of fact.

\emph{Egypt} is the case the story is now being told about, and it is a bid. Bloomberg reported on 26~August, from anonymous sources, that Huawei had bid to build Egyptian government AI datacentres---about 1{,}408 Ascend 950-series accelerators for a training cloud and roughly 600 for two inference clusters, a twelve-month build, reportedly with iFlytek as partner---which, if it closed, would be the first delivered Ascend installation abroad (Brazil's Rio machine is contracted but not due until July 2027)~\cite{huawei-egypt} [A]. The US State Department responded by contacting NVIDIA, AMD and Microsoft to assemble a counter-offer and by warning Cairo that Ascend use might breach US rules, Egypt having needed a US licence for advanced AI chips since 2023. No price has been reported, no financing terms have surfaced, and nothing has been signed; Xi Jinping's state visit to Cairo on 31~August--2~September, the obvious occasion for an announcement, produced no announced AI agreement~\cite{huawei-egypt} [A]. \emph{Malaysia} is an evaluation. Telekom Malaysia is reported to be assessing the Ascend 910C as the backbone of a RM2B (about US\$494M) sovereign project whose purpose is to keep government, military and intelligence data in-country; the quantity is unclear, the project is unsigned, and Malaysia accepted export-control commitments in its 2025 reciprocal trade agreement with the United States, whose officials had already warned in 2025 that 910C use could violate US rules~\cite{huawei-malaysia} [A]. The precedent inside the case is the one to weigh: in May 2025 a Malaysian deputy minister announced a 3{,}000-chip Ascend deployment and the government retracted it within hours. \emph{Brazil} is a contract, and it is the only one. On 21~August 2026 Brazil launched an R\$2.3B (about US\$444M) AI supercomputing programme of which R\$1.3B (about US\$251M) goes to Huawei and iFlytek for a Rio de Janeiro machine on Ascend and a Portuguese-language model, from July 2027; a further R\$1B tender for a second machine in Rio Grande do Norte is expected to go to NVIDIA; the whole is financed by Brazil's own FNDCT science fund, with 3{,}000 professionals to be trained over five years; and the stated principle of the programme is ``not to depend on a single company, technology or country''~\cite{brazil-ai-supercomputer} [A]. Two further entries belong in the ledger below the three. Kazakhstan's 16~July package with China---``\$15B, seventy-plus agreements,'' signed while Tokayev was in Shanghai for the WAICO founding---includes a strategic partnership between the Ministry of AI and Digital Development and Huawei, a Samruk-Kazyna agreement to acquire Huawei equipment, and a Kazakhtelecom--Hengtong agreement on ``basic principles'' for a 1\,GW ``Data Center Valley'' at Ekibastuz with a first facility targeted for June 2027; project-level values are undisclosed and the package is a mix of binding agreements and memoranda~\cite{kazakh-huawei} [A]. And Pakistan's Punjab government agreed six datacentres and a thousand ``AI master trainers'' with Huawei on 28~August, with no money, no chips and no contract described~\cite{kazakh-huawei} [A]. Table~\ref{tab:export-offer} lays the five out on the four questions that matter.

\begin{table}[htbp]
  \centering\footnotesize
  \caption{The Chinese AI export offer as it stood at \datafreeze. One contract, one bid, one evaluation, two frameworks; no financing from the vendor's side in any row~\cite{huawei-egypt,huawei-malaysia,brazil-ai-supercomputer,kazakh-huawei}.}
  \label{tab:export-offer}
  \setlength{\tabcolsep}{4pt}
  \begin{tabular}{p{17mm}p{40mm}p{22mm}p{25mm}p{30mm}}
    \toprule
    Country & What was offered & Status & Financing & US response \\
    \midrule
    Egypt & Huawei bid: $\sim$1{,}408 Ascend 950-series for training, $\sim$600 for inference; 12-month build; iFlytek reported & Bid; unsigned; no price; no AI agreement from Xi's Cairo visit & None reported & State Dept assembling NVIDIA/AMD/Microsoft counter-offer; warning that Ascend use may breach US rules \\
    \addlinespace
    Malaysia & Ascend 910C as backbone of a RM2B ($\sim$US\$494M) sovereign project (Telekom Malaysia) & Under evaluation; unsigned; 2025's 3{,}000-chip announcement retracted within hours & Buyer's own budget & 2025 warning on 910C; export-control commitments in Malaysia's 2025 US trade agreement \\
    \addlinespace
    Brazil & R\$1.3B ($\sim$US\$251M) to Huawei/iFlytek: Rio machine on Ascend plus Portuguese-language model, from July 2027 & Contracted (21 Aug.\ 2026); second R\$1B machine expected to go to NVIDIA & Brazil's FNDCT; no Chinese credit & None reported \\
    \addlinespace
    Kazakhstan & Huawei strategic partnership; Samruk-Kazyna equipment purchase; 1\,GW ``Data Center Valley'' on ``basic principles'' & MoUs and agreements; values undisclosed; first facility targeted June 2027 & Undisclosed & None reported \\
    \addlinespace
    Pakistan (Punjab) & Six datacentres; 1{,}000 ``AI master trainers'' & Agreement (28 Aug.\ 2026); no chips, money or contract described & None described & None reported \\
    \bottomrule
  \end{tabular}
\end{table}

\subsection{The two offers, side by side}

The American offer must be held to the same standard, and it fares no better. Executive Order 14320, ``Promoting the Export of the American AI Technology Stack'' (23~July 2025), and the Commerce Department's American AI Exports Program that implements it describe exactly what the Chinese wrapper does not: full-stack packages---hardware, data, models, security, applications---with federal financing referrals to direct loans, guarantees, equity and political-risk insurance, and an Export--Import Bank ``ExportAI'' framework under consideration since May 2026~\cite{us-ai-exports} [P]. Seventy-eight consortium proposals had been received by 30~June 2026. No package had been designated at \datafreeze~\cite{us-ai-exports} [P]. So the ledger reads: the United States has a programme, a financing architecture on paper, seventy-eight bids and zero designations; China has no comparable programme, no financing facility, no designated packages because there is nothing to designate, and no confirmed Ascend export. \emph{Neither offer has been delivered anywhere.} The contest in Cairo that the reporting describes as a proxy war between two stacks is, precisely stated, a contest between a bid and a counter-offer, neither of which exists in the field, for a buyer who has signed nothing with either~\cite{huawei-egypt,us-ai-exports}. The buyer that \emph{has} signed---Brazil---signed with both, financed itself, and said in terms that it was doing so in order not to depend on either.

\subsection{Where the telecom template transfers, and where it does not}

Chapter~\ref{ch:telecom} supplied the template that the sovereign-pact story is borrowed from, and closed by naming three of its elements that do not transfer; this section is where each is tested. The instrument that built Huawei's network position was credit: a China Development Bank facility of US\$10B for Huawei's customers in 2004, raised to US\$30B in December 2009, on which an operator in Lagos or Jakarta financed a switch it could not otherwise have bought. The condition that made the credit cheap was surplus: factories the state wanted running. And the object that held the customer was a bundle of hardware, service and debt that bound the operator for a generation. Test each against the AI record and each fails, and the failures are of different kinds.

\emph{The credit is absent.} No China Development Bank or Export--Import Bank facility for AI infrastructure exports had surfaced at \datafreeze. This is an absence in the record, not a verified absence---a facility could exist undisclosed, and Chinese policy finance is not transparent---but the one place where such a facility would have shown itself is the one contract that exists, and Brazil financed that contract from its own science fund. \emph{The capacity runs the other way.} Huawei exported switches from surplus; it cannot export accelerators from surplus because there is none. DeepSeek's reported order of at least 160{,}000 Ascend 950DT for a single inference site is expected to take up to a year to fill, because delivery is paced by Huawei's and the domestic HBM chain's output~\cite{deepseek-ascend-order} [A]; a SemiAnalysis-derived estimate puts 2026 domestic HBM at enough for some 250{,}000--300{,}000 Ascend 910C-equivalents [A/C]; and Beijing's reported requirement that state-funded compute run at least 80\% domestic silicon under a roughly RMB~2T national programme pulls every chip the fleet produces toward the home market before any can be spared for Cairo~\cite{ascend-supply-ceiling} [C]. The same source records that Huawei was pitching low thousands of 910B to the UAE, Saudi Arabia and Thailand in July 2025 with no finalized deals, and that the Gulf went to American suppliers; and that BIS guidance since May 2025 holds that using Ascend ``anywhere in the world'' risks violating US controls, which is the legal instrument behind the warnings to Cairo and Kuala Lumpur~\cite{ascend-supply-ceiling} [P/A]. A vendor whose largest domestic customer absorbs a year of output, whose government reserves the output for domestic use, and whose product is a controlled item in the recipient's jurisdiction is not running the 2009 playbook, whatever it intends. \emph{And the artifact does not lock in.} The diffusion mechanism that actually operates in the non-aligned world is not a datacentre; it is a download. Hugging Face counts 151{,}448 derivative models of the Qwen family, about 2.6 times Meta's, and roughly 39.6 million monthly GGUF downloads for Qwen against 7.5 million for Llama---something over five times the volume~\cite{qwen-derivatives} [V]. Those weights cost the recipient nothing and bind the recipient to nothing; the same buyer that will not sign for Ascend runs Qwen already. That is a better diffusion instrument than Huawei ever had in telecommunications and a far worse capture instrument, and it is the instrument Chinese AI is actually being exported with.

The conclusion is stated plainly because the story it corrects is stated plainly everywhere else. At \datafreeze\ the sovereign-pact account of Chinese AI export is a \emph{story}: the telecom template is visible in intent and in the institutional wrapper, absent in credit, absent in capacity, and being run---to the extent anything is being run---with an instrument, open weights, that diffuses better and captures worse than the one the template was built on. Three things cut against that conclusion and are printed with it. Intent precedes capacity in every precedent campaign in Part~I; the CDB facility of 2004 was opened for a company that had been reselling imported switches seventeen years earlier, and the demand-before-supply move of Chapter~\ref{ch:playbook} is by construction a commitment made before the factories exist, so the absence of a credit line in 2026 is evidence about 2026 and not about 2028. The Egypt bid is the first datacentre-scale Ascend export bid to a government on the record, and its outcome---signed on Chinese terms, signed on American terms, or split as Brazil split it---is the test this section will be scored against. And Brazil is not nothing: it is a real contract, on Chinese silicon, for a sovereign model in a G20 country's language, paid for by a G20 government that had every alternative available. One such contract is not a campaign. It is, however, the first row of the table a campaign would be measured in.

\section{Nine positions}

\textbf{Japan} holds the strongest technical hand of any non-principal and the weakest demand-side position. Assets: an exascale HPC lineage whose architectural argument this book carries (Chapters~\ref{ch:compute}, \ref{ch:sizing}); the most defensible remaining semiconductor materials and equipment franchises; leading-edge memory and packaging supply chains; the only measurement-anchored national AI4S demand model in existence; and unusual standing in both blocs' technical communities. The 2025--26 record shows the machinery engaging: the AI Promotion Act (May 2025)---Japan's first AI-dedicated law, deliberately innovation-first, penalty-free, and agile-governance-shaped, the philosophical antipode of the AI Act~\cite{japan-ai-act}; FugakuNEXT contracted to Fujitsu and NVIDIA (Aug.\ 2025) as a MONAKA-X Arm CPU tightly coupled to GPUs over NVLink Fusion, with a $\sim$¥110B initial budget toward operations around 2030~\cite{fugakunext}; the GENIAC compute-subsidy program in its fourth cycle; Sakana at a \$2.65B valuation; Rapidus running a 2\,nm pilot line toward 2027 mass production; SoftBank converting a Sharp LCD plant into a 150--240\,MW AI data center as the anchor of a domestic OpenAI venture; and a fiscal commitment moving from the ¥10T-through-2030 package toward a quadrupled annual chips-and-AI budget~\cite{japan-budget}. Liabilities, unchanged and now quantified: a domestic model ecosystem positioned for Japanese-language enterprise depth rather than the open frontier, an energy position that cannot win a watts race, and a research population whose committed provisioning sits an order of magnitude below its own modeled requirement. The honest tension the record exposes: FugakuNEXT's architecture buys NVIDIA coupling at the exact moment this book's compute chapter argues the accelerated-CPU corner is the sovereign long game---defensible as a hedge, dangerous as a habit. Recommended posture: build the accelerated-CPU line to its FP8-capable conclusion \emph{alongside} the GPU coupling, not instead of sovereignty; adopt the best open weights irrespective of origin under a full trust stack; bid for the certification-authority vacancy in coalition rather than alone; and use the sizing framework as the diplomatic instrument it is---a shared unit is a claim on the terms of the debate.

\textbf{The European Union} has the pen and not the press---and in 2025--26 the pen visibly hesitated. The Digital Omnibus process pushed the AI Act's high-risk obligations back by sixteen months to two years (Annex III to December 2027, embedded systems to August 2028) while strengthening the AI Office's GPAI powers---an implicit concession that the regulatory clock had outrun the industrial one~\cite{eu-omnibus}. Meanwhile the industrial program scaled: thirteen AI factories plus antennas under a $\sim$€10B envelope, and a formal EuroHPC tender (July 2026) for up to seven gigafactories of 100k-plus accelerators each, against 76 expressions of interest totaling $\sim$\$230B---an oversubscription that says European capital wants sovereign compute even where European regulation slows its use~\cite{eu-gigafactories,eu-giga-tender}. Mistral closed €1.7B at $\sim$\$14B with \emph{ASML} as lead investor---the continent's one true chokepoint firm buying into its one frontier-adjacent lab, sovereignty consolidating by cap table~\cite{mistral-asml}---and JUPITER delivered Europe's first exascale machine. Liabilities: fragmented procurement, energy costs, no top-tier frontier laboratory, and the persistent gap between regulatory ambition and industrial capacity that the omnibus retreat priced publicly. Recommended posture, sharpened by the record: spend regulatory energy on the \emph{conformity infrastructure} the world lacks---evaluation batteries, signed provenance, notified-body capacity for model artifacts, the one arena where the AI Act's machinery is an asset rather than a tax; buy the open commons and audit it; and treat the gigafactory programme as a sovereignty asset only if paired with the operating capability and the model programs to use it---seven empty cathedrals of compute would be the solar-tariff mistake with better architecture.

\textbf{India} is the largest genuinely undecided market and the one whose choice would most move the aggregate---and its 2025--26 record is the world's cleanest experiment in demand-side sovereignty. The IndiaAI Mission empanelled $\sim$38{,}000 GPUs across three tenders at subsidized rates reportedly under \$1/GPU-hour---among the cheapest sovereign compute on earth---and then discovered the harder half of the problem: reported utilization lagged badly, a national GPU pool in search of users~\cite{indiaai}. The lesson generalizes to every country in this chapter and confirms this book's cost-stack analysis: \emph{compute is the easy term; $c_{\mathrm{int}}$ and $c_{\mathrm{org}}$ are the binding ones}, and a subsidy that buys hardware without buying integration capacity buys idle silicon. The sovereign-model program (Sarvam on 4{,}096 subsidized H100s, BharatGen, Krutrim's frugal-AI line) and \$45B-plus of announced hyperscaler investment (Microsoft's \$17.5B, Google's \$15B hub, AWS) complete the picture: both blocs building the physical layer, India supplying the population. Recommended posture: aggressive adoption of open weights with domestic fine-tuning and Indic-language specialization; national evaluation capacity; fix the utilization gap before expanding the pool---and the highest-leverage move available remains becoming the world's integration layer, since whoever performs the integration accumulates the complementary capital that the general-purpose-technology literature says captures the value.

\textbf{South Korea} deserves its own position rather than a footnote, because it is the only non-principal that is simultaneously a critical \emph{supplier} to one bloc and a determined sovereign builder. The supplier position is decisive: SK hynix holds $\sim$62\% of HBM and reportedly two-thirds of the leading HBM4 volume---Korea sits on the choke half of the chokepoint this book's compute chapter is about~\cite{korea-hbm}. The sovereign program is the most structured in the world: a national AI computing center, a 260{,}000-GPU national deal spanning government and the chaebol ``AI factories,'' a $\sim$100T-won public-private pledge---and, uniquely, a sovereign foundation-model \emph{tournament}: five teams selected with compute and data support, then eliminated in public rounds, with 700B-parameter open models released as the competition sharpened~\cite{korea-sovereign}. That is playbook step~4 (selection, not planning) executed by a democracy, and this book registers it as the most interesting institutional experiment in national AI policy anywhere: state-funded competition with open-weight release as the qualifying bar. Recommended posture: keep the tournament honest (the temptation to anoint a national champion early is playbook step~6 (cartelize and gatekeep) arriving prematurely); price the HBM franchise as the strategic rent it is while it lasts; and hedge the memory cliff, because Korea is the country most exposed to exactly the capacity-first arbitrage of Chapter~\ref{ch:compute}.

\textbf{Singapore and ASEAN} illustrate the small-state problem in its purest form. Singapore's entire national AI4S requirement is smaller than one large country's first procurement stage (Chapter~\ref{ch:sizing}), which makes autonomy through scale impossible and autonomy through \emph{selection} essential: pick workloads where a small, excellent, well-instrumented capability beats a large generic one, and pool everything else---SEA-LION, the region's open multilingual model family, is the template: small, targeted, open, and adopted because it serves eleven languages the frontier ignores~\cite{sea-lion}. ASEAN more broadly is the clearest case of the third-market dynamic---price-sensitive, sovereignty-conscious, and courted by both blocs with training, infrastructure and, on the American side, the promise of financing, with \$15B-plus of data-center investment landing in Johor alone. It is also where Chinese adoption has moved fastest, for the same reasons Chinese EVs did---and where the limits of that adoption at the hardware layer are most visible: Malaysia's Ascend evaluation is unsigned, its 2025 predecessor was retracted within hours, and its 2025 trade agreement with Washington carries export-control commitments that a signature would test (Section~\ref{sec:export-offer})~\cite{huawei-malaysia} [A].

\textbf{The Gulf states} are the anomaly: capital-rich, energy-rich, compute-poor by choice rather than constraint, and buying position in both stacks simultaneously at a speed neither principal fully controls. The 2025--26 record removed any doubt about scale: Saudi HUMAIN building toward 6.6\,GW with NVIDIA ``several hundred thousand GPUs'' and an AMD \$10B/6\,GW joint venture; Stargate UAE's first gigawatt inside a planned 5\,GW Abu Dhabi campus; MGX targeting \$100B of AI assets as anchor investor in the American frontier itself; and---the geopolitical hinge---US Commerce authorizing 70{,}000 advanced-chip exports to G42 and HUMAIN under security conditions, making the Gulf the first place Washington has \emph{chosen} to locate frontier-scale capacity it controls only contractually~\cite{gulf-ai}. They matter to this book's argument for one specific reason: they are the only non-aligned actors who can bid for frontier-scale capacity on the open market, which makes them the swing buyers in any repricing scenario (Chapter~\ref{ch:trade}) and the most likely hosts of the first sovereign frontier training run outside the two principals---now with the silicon on order.

\textbf{Taiwan} is the position everyone else's position depends on and the hardest one to hold: near-monopoly production of the advanced logic both blocs' AI machines require, a \$3.2B ``ten major AI infrastructure'' program of its own---and an energy system, post-nuclear phase-out, whose arithmetic is the binding constraint on both its sovereign-AI ambitions and its foundry expansion~\cite{taiwan-energy}. Taiwan is this book's busbar chapter compressed into one island: the world's most valuable compute, gated by the world's tightest watt. Recommended posture: treat energy as the semiconductor policy it is; and hold the TSMC franchise as the alliance instrument it has become, while building the sovereign-model capacity (TAIDE-line) that keeps dependence from becoming identity.

\textbf{Brazil and Latin America} combine large populations, real scientific communities, severe fiscal constraints, and---in Brazil's case---a clean grid that is an underrated asset in an energy-constrained industry. The rational posture is almost pure commons adoption plus regional pooling, and the risk is the classic one: importing a service, building nothing, and discovering in a decade that the standards were set elsewhere. Brazil is, at \datafreeze, the one government outside China to have executed the posture in hardware rather than in a strategy document, and the shape of what it executed is the instructive part: a programme split between an Ascend machine with a Portuguese-language model and an NVIDIA machine, financed from its own science fund rather than from either principal's credit, with three thousand people to be trained on it (Section~\ref{sec:export-offer})~\cite{brazil-ai-supercomputer} [A]. That is the six moves below, signed; it is also a single contract of a quarter-billion dollars, and the standards risk it hedges against is not answered by hedging the vendor.

\textbf{Africa} is treated in most Western AI strategy documents as a recipient and in Chinese AI diplomacy as a partner, and the difference in framing is itself strategic. The first facts on the ground arrived in 2026: Cassava's NVIDIA-powered ``AI factory'' in South Africa---the continent's first---with Kenya, Nigeria, Egypt, and Morocco expansions planned, pitched explicitly as sovereign data and compute~\cite{cassava}; and small open models for African languages (InkubaLM-class) following the SEA-LION template. The continent's relevant facts are a young technical population, the fastest-growing developer base in the world, minimal legacy infrastructure to defend, and acute exposure on energy and connectivity. Chapter~\ref{ch:consortium} argues the correct treatment is as a source of talent, scientific questions, and data---not aid---and notes that the institution which builds a real access track will acquire an asset that compounds for decades. Nobody has built it. The World AI Cooperation Organization is bidding---with, so far, a training pledge of five thousand places and no compute, financing or datacentre clause in any of its founding texts (Section~\ref{sec:export-offer}); the one African government reported to be negotiating Chinese silicon, Egypt, is not a WAICO signatory~\cite{waico,huawei-egypt} [A].

\section{What the matrix predicts}

Apply Table~\ref{tab:decision} honestly and an uncomfortable pattern emerges. On \emph{cost}, \emph{sovereignty} (operational), and \emph{export-control exposure}, the Chinese open stack currently scores better for almost every non-aligned actor---it is cheaper per task, it can be run on-premises without foreign permission, and no government can revoke a weight file already downloaded. On \emph{capability at the absolute frontier}, \emph{trust}, and \emph{interoperability with existing enterprise systems}, the American stack scores better. On \emph{standards} both create dependence, in different currencies. And on \emph{industrial spillover}---the dimension that determines whether a country is building anything---\emph{neither stack helps unless the adopting country does the work itself}, which is the single most important sentence in this chapter.

That pattern was drawn before the second commons existed, and Section~\ref{sec:two-commons} moves three of its cells without overturning it. On \emph{cost} the Chinese stack no longer holds an exclusive advantage, because a free artifact is free in either language; what remains of the difference is the hardware the artifact prefers, which is where it always was. On \emph{operational sovereignty} the two are now close to level---a self-hosted Nemotron runs without foreign permission exactly as a self-hosted Qwen does---and the asymmetry survives only at the layer beneath, where one stack's preferred accelerator is a controlled item and the other's is designed not to be. On \emph{export-control exposure} the gap therefore narrows at the model layer and persists intact at the compute layer, which sharpens the finding rather than reversing it. What has not moved is the row that decides the rest: on \emph{trust}, neither commons supplies independent assurance about itself, and a vendor-sponsored artifact nobody has evaluated is exactly as unevaluated as a state-adjacent one. The vacancy identified above as the Chinese position's strategic vulnerability is now a vacancy in both stacks at once---larger, more valuable, and harder for either principal to fill credibly, because in both cases the party best placed to certify the artifact is the party that paid for it.

That pattern explains observed behaviour better than allegiance does. It explains why institutions that ban a Chinese chat application nonetheless evaluate its weights for self-hosted deployment; why sovereignty-conscious governments buy American cloud and Chinese weights simultaneously; and why the fastest-growing adoption of Chinese open models is in exactly the countries with the least ability to pay rent. It also explains a strategic vulnerability of the Chinese position that follows directly from the same logic: the commons currently wins on cost and operational sovereignty but loses on trust, and trust is the dimension that a coalition of non-aligned states could supply \emph{for themselves}---which would make the commons safe to adopt and simultaneously deprive its sponsor of the dependence it purchased.

\section{The non-aligned strategy, stated once}

For any country in this chapter, six moves follow from the matrix, and they are the same six regardless of size:

\emph{Measure your own demand} before procuring anything, in a unit others recognize (Chapter~\ref{ch:sizing}); a country that cannot state its requirement cannot audit a vendor or join a federation. \emph{Adopt the commons---now plural---and audit both its behaviour and its provenance}: the cost and sovereignty arguments are decisive, the trust gap is closable with infrastructure rather than with prohibition, and the provenance questions of Section~\ref{sec:two-commons} should be put to every artifact from either commons, because the licence file is the least informative document in the procurement pack. \emph{Mirror what you depend on}---weights, datasets and evaluation harnesses, in a repository you operate---because every continuity risk in Table~\ref{tab:two-commons} is neutralized by a copy, and a copy is the cheapest sovereignty instrument in this book. \emph{Retain a pretraining capability} sized to competence rather than to leaderboards, because the return is people and the alternative is permanent tenancy. \emph{Pool internationally} for the capacity and the certification you cannot afford alone (Chapter~\ref{ch:consortium}). And \emph{do the integration work domestically}, because the general-purpose-technology literature and this book's Chapter~\ref{ch:theory} agree that complementary capital is where diffused value is captured, and it is the only part of the stack no foreign supplier can perform for you.

The countries that moved early in solar, batteries, and EVs are the ones the record remembers kindly. The window in AI is measured in quarters rather than decades, because the standards layer is being set now and the certification layer is still vacant.
